\section{Local cite-lifts and residual packaging advancements}\label{sec:localadv}
This section records three Category~A cite-lifts from the residual path-problem (RPP) group-ring calculus---named \emph{Local cite-lifts F1--F3} to distinguish them from the unrelated $E_8$ filter families F1--F3 of Section~\ref{sec:emit}---together with the November~2026 packaging advancements that fix the ambient Hilbert newspace, the thin form complex, and the static cell budget. The Local cite-lifts are independent of the local modular constituents of Section~\ref{sec:modlocal}. Nothing in this section settles the residual question for the Bridge Series~$C$, identifies the shadow~$\Xi$, or produces an explicit formula for a boundary vector of~$C$ inside~$\VV$.

Throughout, fix an integer $M\ge2$ and write $G=\Z/M\Z$. Let $\Z[G]$ be the group ring, and let $\rho\colon\Z[G]\to\mathrm{Mat}_M(\Z)$ be the regular representation, so that for $h=\sum_g a_g[g]$ the matrix $\rho(h)$ is the circulant with eigenvalues $\hat h(k)=\sum_g a_g\,\omega^{kg}$, $\omega=e^{2\pi i/M}$, and $\det\rho(h)=\prod_k\hat h(k)$. Given a step alphabet $C\subset\Z$, weights $w=(w_1,\dots,w_n)$ and integer multiplicities $a_{i,c}$, set
\[
h_i:=\sum_{c\in C}a_{i,c}\,[c\cdot w_i]\in\Z[G],\qquad
T:=h_n\cdots h_1,\qquad
|T|:=|h_n|\cdots|h_1|,
\]
where $|h|=\sum_g|a_g|[g]$. For a source residue $r_s$ and a target residue $r_t$, write $W=W(r_s\to r_t)$ for the signed path weight $(\rho(T))_{r_t,r_s}$ and $N=N(r_s\to r_t)$ for the unsigned path weight $(\rho(|T|))_{r_t,r_s}$. Write $\|\cdot\|_1$ for the $\ell^1$-norm of coefficients and $\mathrm{spr}$ for the spectral radius over $\C$.

\subsection{Local cite-lift F1: weak NP-completeness on an odd-$M$ slice}
\begin{theorem}[Local cite-lift F1]\label{thm:local-F1}
Consider the residual path problem (RPP) of deciding whether there exists a path of signed weight summing to a prescribed target residue modulo~$M$. Restrict to the slice in which $M$ is part of the input, $C=\{0,1\}$, every multiplicity equals~$1$, and
\[
M=2^k+1>\Sigma:=\sum_{i=1}^n w_i
\]
for some $k\ge1$. Then:
\begin{enumerate}[label=(\roman*)]
\item for every odd $M$, one has $1+\omega^j\neq0$ for all $j$, hence $\det\rho(h_i)\neq0$ for each layer and $\det\rho(T)\neq0$;
\item RPP restricted to this slice is weakly NP-complete (membership via a polynomial certificate; hardness by reduction from Subset Sum through the odd modulus $M=2^k+1>\Sigma$).
\end{enumerate}
The same decision problem is false for fixed odd $M\le\Sigma$: the instance $M=3$, $w=\{1,2\}$, target~$0$ has no nonempty subset summing to~$0$ over~$\Z$, yet $\{1,2\}\equiv0\pmod3$. Completeness is weak only: a pseudo-polynomial dynamic program of cost $O(nM)$ applies.
\end{theorem}
\begin{proof}
(i)~If $1+\omega^j=0$ then $\omega^j=-1$, so $(-1)^M=\omega^{jM}=1$, contradicting $M$ odd. Hence each $\hat h_i(k)=1+\omega^{k w_i}$ (or $2$ when $w_i\equiv0$) is nonzero, and $\det\rho(T)=\prod_i\det\rho(h_i)\neq0$. (ii)~The standard Subset-Sum reduction with modulus $2^k$ produces layers with $\det\rho(h_i)=0$ whenever some $w_i$ is odd (evaluate at frequency $M/2$). Replacing the modulus by the odd integer $M=2^k+1>\Sigma$ restores nonvanishing of every determinant while preserving the subset-sum encoding of yes/no instances, because the encoding uses only the inequality $M>\Sigma$ (and, for no-instances, a reachable-residue argument independent of parity). Membership is the usual short certificate that a path realises the target; checking that the instance lies on the slice is polynomial for $C=\{0,1\}$ and unit coefficients. The counterexample at fixed $M=3$, $w=\{1,2\}$, target~$0$ is immediate. Weakness follows from the standard $O(nM)$ DP on residues.
\end{proof}
\catA{Cite-lift. Weak NP-completeness of RPP on the odd-$M=2^k+1>\Sigma$ unit-coefficient slice with $C=\{0,1\}$; false for fixed odd $M\le\Sigma$. Distinct from the $E_8$ filter family F1 of Section~\ref{sec:emit}.}

\subsection{Local cite-lift F2: nonvanishing determinant does not prevent cancellation}
\begin{proposition}[Local cite-lift F2]\label{prop:local-F2}
Nonvanishing of $\det\rho(h_i)$ does not preclude signed path cancellation. Explicitly, for $M=3$,
\[
h_1=[0]+[1],\qquad h_2=-[0]+[1]+3[2],
\]
one has $\det\rho(h_1)=2$ and $\det\rho(h_2)=36$, yet at target residue~$1$ the two contributing paths cancel: $W=0$ and $N=2$.
\end{proposition}
\begin{proof}
The eigenvalues of a circulant are the DFT evaluations, so $\det\rho(h_1)=\prod_{k=0}^{2}(1+\omega^k)=2$ and $\det\rho(h_2)=\prod_{k=0}^{2}(-1+\omega^k+3\omega^{2k})=36$. The two length-$2$ paths from residue~$0$ to residue~$1$ have signed weights $1\cdot(-1)$ and $1\cdot1$, summing to $W=0$, while the unsigned count is $N=2$.
\end{proof}
\catA{Cite-lift. Explicit $M=3$ counterexample: $\det\rho\neq0$ coexists with $W=0$, $N=2$. Distinct from the $E_8$ filter family F2 of Section~\ref{sec:emit}.}

\subsection{Local cite-lift F3: signed bound by the regular spectral radius}
\begin{proposition}[Local cite-lift F3]\label{prop:local-F3}
For arbitrary integer multiplicities,
\[
|W|\;\le\; N\;\le\;\prod_i\|h_i\|_1\;=\;\mathrm{spr}\bigl(\rho(|T|)\bigr).
\]
In particular the inequality requires no nonvanishing-determinant hypothesis. The stronger reading that nonvanishing of $\det\rho$ \emph{causes} the bound is false (Proposition~\ref{prop:local-F2}). No claim of ``total cancellation'' is made.
\end{proposition}
\begin{proof}
The triangle inequality gives $|W|=|\sum\sigma|\le\sum|\sigma|=N$. The quantity $N$ is an entry of the nonnegative circulant $\rho(|T|)$; the augmentation $\varepsilon\colon\Z[G]\to\Z$ is a ring homomorphism, so the column sums of $\rho(|T|)$ equal $\||T|\|_1=\prod_i\|h_i\|_1$, and every entry is at most the column sum. For a nonnegative circulant the all-ones vector is a Perron eigenvector with eigenvalue equal to the row sum, hence $\mathrm{spr}(\rho(|T|))=\prod_i\|h_i\|_1$.
\end{proof}
\catA{Cite-lift. The bound $|W|\le N\le\prod\|h_i\|_1=\mathrm{spr}(\rho(|T|))$ holds with no determinant hypothesis; det-causation is excluded. Distinct from the $E_8$ filter family F3 of Section~\ref{sec:emit}.}

\begin{remark}[Scope of the Local cite-lifts]\label{rem:local-scope}
The three Local cite-lifts F1--F3 exhaust the Category~A content drawn from the group-ring reading of the residual path problem. They are combinatorial and scale-free. No normalisation scalar enters them, and they do not identify any modular form, any shadow series, or any boundary vector of~$C$.
\end{remark}

\subsection{Absence of a normalisation for the Local cite-lifts}
\begin{proposition}[Scale-free Local cite-lifts; Hecke-scale fallback]\label{prop:Lambda0-none}
No normalisation scalar enters Theorems~\ref{thm:local-F1}--\ref{prop:local-F3}. In particular the Local cite-lifts are independent of every candidate of the form
\[
\frac{\ln(6+\sqrt{35})}{\sqrt{10}},\qquad
1.81,\qquad
0.7858.
\]
If a later Hecke deformation over $K=\Q(\sqrt{35})$ requires a field-invariant scale, the only canonical choice recorded here is the residue of the Dedekind zeta at $s=1$,
\begin{equation}\label{eq:res-zetaK}
\Res_{s=1}\zeta_K
\;=\;
\frac{2\,h_K\,R_K}{\sqrt{140}}
\;=\;
\frac{2\ln(6+\sqrt{35})}{\sqrt{35}}
\;\approx\;0.837679,
\end{equation}
where $h_K=2$ and $R_K=\ln(6+\sqrt{35})$ (fundamental unit of norm~$+1$). Identity~\eqref{eq:res-zetaK} equals $L(1,\chi_{140})$ and is \emph{not} a normalisation of the Local cite-lifts.
\end{proposition}
\begin{proof}
The statements of Theorems~\ref{thm:local-F1}--\ref{prop:local-F3} are purely combinatorial identities and complexity statements in $\Z[G]$; no scalar appears. The classical formula for $\Res_{s=1}\zeta_K$ over a real quadratic field of discriminant~$140$ is standard ($r_1=2$, $w=2$), and the numerical evaluation follows from $h_K=2$ and $R_K=\ln\varepsilon_K$ with $\varepsilon_K=6+\sqrt{35}$.
\end{proof}
\catA{No normalisation for Local F1--F3; $\Res\zeta_K\approx0.837679$ reserved only as a Hecke-scale fallback.}

\subsection{The full level-one Hilbert newspace over $\Q(\sqrt{35})$}
Write $K=\Q(\sqrt{35})$, so that the LMFDB label of $K$ is $2.2.140.1$. Let $\mathrm{Cl}^+(K)$ denote the narrow class group; one has $h_K=2$ and $h^+=4$, with
\[
\mathrm{Cl}^+(K)\;\cong\; C_2\times C_2
\]
generated by the four narrow classes
\begin{equation}\label{eq:clplus}
\begin{aligned}
C_{00}&=(1),&
C_{10}&=\mathfrak p_2=(2,1+\sqrt{35}),\\
C_{01}&=\mathfrak p_{19},&
C_{11}&=\mathfrak p_5=(5,\sqrt{35}).
\end{aligned}
\end{equation}

\begin{definition}[The space $\mathcal V_{35}^{\mathrm{Hilb}}$]\label{def:V35hilb}
Let $\mathcal V_{35}^{\mathrm{Hilb}}$ denote the full cuspidal newspace of parallel weight $[2,2]$ and level~$(1)$ for Hilbert modular forms over~$K$. As an LMFDB newspace it is the direct sum of the four newforms
\begin{equation}\label{eq:V35split}
\begin{array}{c|ccc}
\text{label} & \dim & \mathrm{CM} & \text{base change}\\\hline
2.2.140.1\text{-}1.1\text{-}a & 2 & \mathrm{yes} & \mathrm{yes}\\
2.2.140.1\text{-}1.1\text{-}b & 2 & \mathrm{yes} & \mathrm{no}\\
2.2.140.1\text{-}1.1\text{-}c & 2 & \mathrm{yes} & \mathrm{yes}\\
2.2.140.1\text{-}1.1\text{-}d & 8 & \mathrm{no} & \mathrm{yes}
\end{array}
\end{equation}
of total dimension $2+2+2+8=14$. The numeral ``$35$'' in the subscript names the field~$K$, not the dimension. The space includes all four narrow class components~\eqref{eq:clplus}.
\end{definition}

\begin{theorem}[Dimension and narrow components]\label{thm:V35}
One has $\dim\mathcal V_{35}^{\mathrm{Hilb}}=14$, with the split~\eqref{eq:V35split} verified against the LMFDB newform pages for $2.2.140.1$ at parallel weight $[2,2]$ and level~$1$. The space contains every narrow class component of $\mathrm{Cl}^+(K)$. These forms are Hilbert modular forms over~$K$, not Siegel modular forms.
\end{theorem}
\begin{proof}
The four newform pages of LMFDB for field $2.2.140.1$, level~$1$, weight~$[2,2]$ list dimensions $2,2,2,8$ with the CM and base-change attributes of~\eqref{eq:V35split}; the sum is~$14$. The narrow class number $h^+=4$ is the standard narrow-class computation for discriminant~$140$, and Definition~\ref{def:V35hilb} includes every component by fiat. Hilbert versus Siegel is the standard dictionary for parallel weight $[2,2]$ over a real quadratic field.
\end{proof}
\catA{Full level-$1$ parallel-weight $[2,2]$ cuspidal newspace over $\Q(\sqrt{35})$, dimension $14=2+2+2+8$, all four $\mathrm{Cl}^+$ components; Hilbert, not Siegel.}

\begin{remark}[Notation relative to $\VV$]\label{rem:V35-vs-VV}
The space $\mathcal V_{35}^{\mathrm{Hilb}}$ of Definition~\ref{def:V35hilb} is not to be confused with the $36$-dimensional span $\VV$ of the $35$-component law vector of Part~III (Definition~\ref{def:K}). No identification between them is claimed.
\end{remark}

\subsection{Classical arithmetic corrections}
\begin{proposition}[Classical exponents and conductors]\label{prop:classical-corr}
The following classical facts are recorded for reference; none is promoted as a bound for the Bridge Series~$C$ or for the Local cite-lifts.
\begin{enumerate}[label=(\roman*)]
\item The classical Burgess exponent on character sums is $3/16$, not $1/5$.
\item The conductor of the elliptic curve $Y^2=X^3-35X$ is $78400$, not $560$. This curve is a different object from $\mathcal V_{35}^{\mathrm{Hilb}}$.
\item The Kim--Sarnak bound $\theta=7/64$ applies to Maass forms; for holomorphic cusp forms of weight~$2$ one has $\theta=0$ by Deligne.
\end{enumerate}
No Burgess bound, no Kim--Sarnak bound, and no conductor-$560$ statement is used as a hypothesis anywhere in this paper.
\end{proposition}
\begin{proof}
(i)~is Burgess's theorem. (ii)~is the Cremona/LMFDB conductor of the curve with Weierstrass model $y^2=x^3-35x$ (minimal discriminant and conductor $78400=2^6\cdot5^2\cdot7^2$). (iii)~is Deligne's theorem for holomorphic weight~$2$, contrasted with the Kim--Sarnak hybrid bound for Maass forms.
\end{proof}
\catA{Recorded classical facts only: Burgess $3/16$; conductor $78400$; Kim--Sarnak $7/64$ not for holomorphic weight~$2$.}

\subsection{The thin form complex and the rank-one classifier $D_1$}
Let $X_\times=K_9\times P_{B'}$ be the product of the complete $9$-vertex charge skeleton with the residual path of length $B'=539$, and let $F$ be the thin residual form complex on the locked generators $\{1,\alpha,\omega_2,\mu,f,\delta f\}$ with product differential~$D$, under the residual relations $\omega_2=d_Q\alpha$ and $\mu=-B'\,\omega_2$ on ordered triples, truncated by $F^3:=0$ (brutal truncation $\sigma_{\le2}$). Write $s_{\alpha\delta}$ for the mixed generator $\alpha\otimes\delta f$.

\begin{theorem}[Rank-one classifier $D_1$]\label{thm:D1}
In the thin truncated complex~$F$ one has
\[
H^0(F)\;\cong\;\Q,\qquad
H^1(F)\;=\;0,\qquad
H^2(F)\;\cong\;\Q\cdot[s_{\alpha\delta}],
\]
and consequently the formality equivalence $F\simeq\Q[0]\oplus\Q[-2]$. Define $D_1$ to be the $\Q$-linear functional on $H^2(F)$ that isolates the line $\Q\cdot[s_{\alpha\delta}]$. Then:
\begin{enumerate}[label=(\roman*)]
\item $D_1$ is a static rank-$1$ classifier over~$\Q$, spanned by $[s_{\alpha\delta}]\in H^2(F)$;
\item $D_1$ acts only on the nine charge sectors of $X_\times$ and takes no argument from any temporal variable, execution schedule, or metric normalisation;
\item no algebraic datum of $D_1$ is shared with any resonant counter set;
\item Layer~$2$ structure beyond the definition of $D_1$ is not constructed from $D_1$ in this paper.
\end{enumerate}
The class $[s_{\alpha\delta}]$ is the generator of $\mathrm{coker}(D^1\colon F^1\to F^2)$ in the thin truncation; in the filled cochain complex of $X_\times$ the same cochain is not a cocycle whenever $\delta f\not\equiv0$.
\end{theorem}
\begin{proof}
The cohomology computation $H^0(F)\cong\Q$, $H^1(F)=0$, $H^2(F)\cong\Q\cdot[s_{\alpha\delta}]$ is the locked computation on the thin residual form complex: $Df=v_\delta$ forces $f\in F^0$, the relation $\ker D=\mathrm{span}\{v_\delta\}=\mathrm{im}(D\colon F^0\to F^1)$ kills $H^1$, and after quotienting by the exact pure-charge pieces the mixed class $[s_{\alpha\delta}]$ spans $H^2$. Formality $F\simeq\Q[0]\oplus\Q[-2]$ is the unique formal model with these Betti numbers. Items~(i)--(iii) are the definition of $D_1$ as the projection onto that line, with domain restricted to the static charge sectors and with no structure map into any resonant counter. Item~(iv) records that no further Layer~$2$ construction is derived from $D_1$ here. The filled-complex caveat is the identity $Ds_{\alpha\delta}=-\delta f$ on prisms, which is nonzero on the residual data.
\end{proof}
\catA{Rank-$1$ over $\Q$, spanned by $[s_{\alpha\delta}]\in H^2(F)\cong\Q$; $H^1(F)=0$; formality $F\simeq\Q[0]\oplus\Q[-2]$; Layer~$2$ not constructed from $D_1$.}

\subsection{Static cell budget and Residual versus Resonant}
\begin{proposition}[Static cell budget $S_{\mathrm{eff}}$]\label{prop:Seff}
Let $T_3$ denote the ternary map of Proposition~\ref{prop:T3} and write $539$ for the residual cardinality of Proposition~\ref{prop:T3}. Define the effective packaging structure $S_{\mathrm{eff}}$ by the following constraints.
\begin{enumerate}[label=(\roman*)]
\item The integer $T_3=539$ is a static cell capacity (equivalently: a static arithmetic hash budget). It is not a stopping time, not a convergence rate, and not a runtime execution length forced by residual cohomology.
\item There is no global guarantee that a packaging remainder $\Omega$ satisfies $\Omega\equiv0\pmod{539}$.
\item Dual-rail representations, when mentioned, are hardware-only side-channel countermeasures and play no algebraic role in $S_{\mathrm{eff}}$.
\item The assignment
\[
\mathcal M\colon S_{\mathrm{eff}}\longrightarrow\bigl(\mathrm{Residual}_{539}\times\mathrm{Resonant}_{539}\bigr)
\]
keeps the two factors disjoint: the residual factor is a count of leftover cells, the resonant factor is a forced round count; a count is not a clock.
\end{enumerate}
\end{proposition}
\begin{proof}
Items~(i)--(iv) are definitional constraints on the packaging arithmetic already present in Proposition~\ref{prop:T3} and Remark~\ref{rem:nolink}: the identity $539=\lfloor4859/9\rfloor=18+521$ records a cell count and, separately, a window-plus-body split used as a schedule; neither identity produces a congruence $\Omega\equiv0\pmod{539}$ nor a structure map from residual cohomology into the schedule. Dual-rail hardware is excluded from the algebraic criteria by~(iii).
\end{proof}
\catA{$T_3=539$ is a static cell/hash budget, not a stopping time; no global $\Omega\equiv0\pmod{539}$; dual-rail hardware-only; Residual~$\neq$~Resonant.}

\begin{proposition}[Residual is not Resonant]\label{prop:res-vs-res}
The residual cardinality $539$ (a count of leftover cells per charge sector after packaging) is a different mathematical object from the resonant schedule length $539$ (a forced iteration count of $T_3$). Equality of the two integers is an equality of values, not a structure map of chain complexes, graded vector spaces, or sets. In particular, residual product cohomology does not force the resonant schedule, and the resonant schedule does not redefine residual cohomology.
\end{proposition}
\begin{proof}
By construction of the residual product complex, $B'=539$ is the number of residual path edges in each sector. By construction of the resonant schedule, the same integer is used as a constant-time iteration budget. No chain map, graded map, or set map relating $C^\bullet(X_\times)$, $F$, or $H^\bullet(F)$ to the counter set $\{1,\dots,539\}$ is part of either definition (Theorem~\ref{thm:D1}(iii) and Proposition~\ref{prop:Seff}(iv)).
\end{proof}
\catA{Residual cardinality $\neq$ resonant schedule; equality of integers is not a structure map.}

\subsection{Status of claims not established}
\begin{proposition}[Claims not part of the theory]\label{prop:not-established}
The following are recorded as not established, and are not used as hypotheses in this paper.
\begin{enumerate}[label=(\roman*)]
\item A Petersson inner-product formula for the shadow~$\Xi$ is not used as a proved identity.
\item No Kronecker $11$-element purge and no four-forces (or $SU(3)$ continuum) mapping is part of the theory.
\item The assertion $C\notin\mathcal V_{35}^{\mathrm{Hilb}}$ still lacks an explicit formula for a boundary vector of~$C$; the elliptic curve $Y^2=X^3-35X$ is not a substitute.
\end{enumerate}
\end{proposition}
\begin{proof}
(i)--(iii) are negative inventory statements: no written Petersson identity for~$\Xi$ is cited as proved; continuum and four-forces maps are excluded from every Category~A statement above; and without an explicit formula for a boundary component of~$C$ the membership $C\notin\mathcal V_{35}^{\mathrm{Hilb}}$ cannot be evaluated.
\end{proof}
\catA{Status only: $\Xi$ Petersson not used as proved; Kronecker/four-forces maps excluded; $C\notin\mathcal V_{35}^{\mathrm{Hilb}}$ lacks an explicit formula for~$C$.}

\begin{remark}[Scope relative to Parts~I--III]\label{rem:localadv-scope}
The Local cite-lifts F1--F3 and the packaging advancements of this section do not modify the radial or $S$-law theorems already proved for the Bridge Series~$C$, do not resolve Problem~\ref{op:xi}, and do not identify the shadow~$\Xi$. The seven Category~A local modular theorems of Section~\ref{sec:modlocal} remain unchanged.
\end{remark}

\noindent\textit{Tip untouched. / No CERT.}
